\documentclass[10pt,a4paper]{amsart}
\usepackage[margin=19mm]{geometry}
\usepackage{amsmath,amssymb,mathtools}
\usepackage[scr=rsfs,cal=euler]{mathalfa}
\usepackage{xfrac}
\usepackage[unicode,hidelinks,bookmarks=false]{hyperref}
\newcommand{\ie}{i.e.\ }
\newcommand{\cf}{cf.\ }
\newcommand{\resp}{respectively\ }
\newcommand{\Subsection}{\subsection}
\providecommand{\nobreakdash}{\nobreak\hbox{-}}
\setlength{\emergencystretch}{2em}
\multlinegap0pt
\usepackage[shortlabels]{enumitem}
\usepackage{tikz}
\usepackage{float}
\usepackage{amssymb}
\usepackage{dsfont}
\newtheorem{remark}{Remark}
\newtheorem{theorem}{Theorem}
\newcommand{\bb}[1]{{\mathbb #1}}
\newcommand{\dd}{\text{\rm d}}
\newcommand{\mf}[1]{{\mathfrak #1}}
\newcommand{\msf}[1]{{\mathsf #1}}
\newcommand{\norm}[1]{\lVert#1\rVert}
\title{A functional central limit theorem for the general Brownian motion on the half-line}
\author{Dirk Erhard, Tertuliano Franco, Milton Jara, Eduardo Pimenta}
\date{Source: arXiv:2408.06830v3 (CC BY 4.0)}
\begin{document}
\maketitle

\noindent\emph{Source and licence.} A functional central limit theorem for the general Brownian motion on the half-line. Version arXiv:2408.06830v3 (CC BY 4.0), distributed by arXiv under the Creative Commons Attribution 4.0 International licence, http://creativecommons.org/licenses/by/4.0/, as recorded in the arXiv licence metadata for this exact version and shown on the abstract page https://arxiv.org/abs/2408.06830v3. Full licence notice: \texttt{licenses/arxiv-2408.06830-CC-BY-4.0.txt}.\par\smallskip

\noindent\textit{Editorial scope.} Counted unit: the article's theorem thm21 (source0.tex lines 270--277). The article's complete original proof of it is delivered with it (source0.tex lines 647--678, one \texttt{proof} environment of 32 lines, which the article places away from the statement). No mathematics is written, repaired or re-derived: the statement and the proof are the article's own lines, in the article's own notation, and no retained line is substituted or re-flowed.  Retained uncounted as the source-local context the proof needs: lines 245--265.  Nothing was omitted from any retained span, so the package carries no omission record.  No retained line contains a citation, so the wrapper carries no bibliography and no bibliography entry.  No retained line contains a figure environment, an image-inclusion command or drawing code, so no image is delivered and the package declares no asset dependency.  Editorial formatting only, in addition: an AMS \texttt{amsart} preamble that restates the article's own declarations and macros used by the retained lines, namely the environments \texttt{remark}, \texttt{theorem} on separate counters, together with the standard packages those lines need.  The retained environments print in this document as Remark 1, Theorem 1 in order of appearance, whereas the article prints the same items with the numbering of its own full document.  The complete native source of the article as distributed in the arXiv e-print for this exact version is bundled unchanged as \texttt{sources/163718/source0.tex}.
	Let $\pi_n : \msf{B} \to \msf{B}_n$ be bounded linear operators indexed by $n \in \bb{N}$, which we call \textit{natural projections} for reasons that will become clear in the sequel.
	Let  $\Xi_n : \mf{D}(\msf{L}^2)\subset \msf{B} \to \msf{B}_n$ be a family of linear operators (not necessarily bounded), which we call \textit{the correction operators}. Denote by $\norm{\cdot}_\text{op}$ the operator norm. We further denote
	\begin{equation}\label{eq:A}
		\Phi_n \;=\; \pi_n + \Xi_n\,.
	\end{equation}
	The following hypotheses will be crucial in this work. Here, by $f \in \mf{D}(\msf{L}^2)$ we mean that $f \in \mf{D}(\msf{L})$ and $\msf{L}f \in \mf{D}(\msf{L})$.
	\begin{enumerate}[label=(H\arabic*)]
		\item \label{(H1)} If $f\in \mf{D}(\msf{L})$, then $\Phi_nf\in\mf D(\msf{L_n})$.
		\item \label{(H2)} There exist sequences $s_1(n)\searrow0$, $s_2(n)\searrow0$ and $s_3(n)\searrow 0$ such that, for any $f\in\mf{D}(\msf{L}^2)$,
		\begin{equation*}
			\norm{\pi_n\msf{L}f-\msf{L}_n\Phi_nf}\;\leq\;s_1(n)\norm{f}+s_2(n)\norm{\msf{L}f}+s_3(n)\norm{\msf{L}^{2}f}\,.
		\end{equation*}
		\item \label{(H3)} There exist sequences $r_1(n)\searrow0$ and $r_2(n)\searrow0$ such that, for each $f\in\mf{D}(\msf{L})$,
		\begin{equation*}
			\norm{\Xi_nf}\;\leq\;r_1(n)\norm{f} +r_2(n)\norm{\msf{L}f}\,.
		\end{equation*}
		\item \label{(H4)} For each $f\in \mf{D}(\msf{L}^2)$, the function $s \in [0,t] \mapsto \Phi_n T(t-s)f\in \msf{B}_n$ is differentiable and its derivative is given by $-\Phi_n \msf{L} T(t-s) f$.
	\end{enumerate}
	\begin{remark}\label{rmk21}
		Since $\pi_n$ is a linear bounded operator, it is easy to check that   $s \in [0,t] \mapsto \pi_n T(t-s)f\in \msf{B}_n$ is differentiable with derivative  $-\pi_n \msf{L}T(t-s) f$. The key point in \ref{(H4)} is to assure that    $s \in [0,t] \mapsto \Xi_n T(t-s)f\in \msf{B}_n$ is differentiable with derivative  $-\Xi_n \msf{L} T(t-s) f$.
	\end{remark}

	\begin{theorem}\label{thm21}
		Under hypotheses  \ref{(H1)} -- \ref{(H4)}, for any $f \in \mf{D}(\msf{L}^2)$ and for each $t$ in a compact interval $[0,b]$ we have that
		\begin{align*}%\label{convaprox}
			\norm{{T}_n(t)\pi_nf - \pi_nT(t)f}\;\lesssim\; & \max\left\{s_1(n),r_1(n)\right\}\norm{f} +\max\big\{s_2(n),r_1(n),r_2(n)\big\}\norm{\msf{L}f}\\
			&+\max\big\{s_3(n),r_2(n)\big\}\norm{\msf{L}^2f}\,.
		\end{align*}
		Here, the proportionality constant depends only on $b$.
	\end{theorem}

	\begin{proof}[Proof of Theorem \ref{thm21}]
		Putting together hypotheses \ref{(H2)} and \ref{(H3)},   we can infer that
		\begin{align}\label{eq3.1}
			\norm{\Phi_n\msf{L}f - \msf{L}_n\Phi_nf}\,&\leq\,s_1(n)\norm{f}+2\max\{s_2(n),r_1(n)\}\norm{\msf{L}f}+2\max\{s_3(n),r_2(n)\}\norm{\msf{L}^2f}\,,
		\end{align}
		for any $f \in \mf{D}(\msf{L}^2)$.
		Fix $t \geq 0$ and define $g_{s,t} := \msf{T}(t - s)f$ for $0 \leq s \leq t$. Then,
		\begin{equation}\label{derivative}
			\partial_s g_{s,t}\;=\;\partial_s\msf{T}(t\,-\,s)f\;=\;-\,\msf{L}\msf{T(t\,-\,s)f\;=\;-\,\msf{L}g_{s,t}}.
		\end{equation}
		Note that $\Phi_nf = \Phi_ng_{t,t}$. From \eqref{derivative}  we obtain that
		\begin{align*}
			\msf{T}_n(t)\Phi_nf\,&=\,\msf{T}_n(0)\Phi_ng_{0,t}\,+\,\int_0^t\partial_s(\msf{T}_n(s)\Phi_ng_{s,t})\,\dd s\\
			&\overset{\ref{(H4)}}{=}\,\msf{T}_n(0)\Phi_ng_{0,t}\,+\,\int_0^t\big[\msf{L}_n\msf{T}_n(s)\Phi_ng_{s,t}\,-\,\msf{T}_n(s)\Phi_n\msf{L}g_{s,t}\big]\,\dd s\,.
		\end{align*}
		Using the equality above, using that semigroups are contractions, that $T_n(0)$ is the identity mapping and using that  $\msf{T}(t)\msf{L} = \msf{L}\msf{T}(t) $ and $\msf{T}_n(t)\msf{L}_n = \msf{L}_n\msf{T}_n(t) $ for all $n \in \bb{N}$ and for all $t \geq 0$,  we have that
		\begin{align}
			\norm{\msf{T}_n(t)\Phi_nf-\Phi_n\msf{T}(t)f}\;=\;&\Big\Vert\int_0^t\msf{T}_n(s)\big[\msf{L}_n\Phi_ng_{s,t}-\Phi_n\msf{L}g_{s,t}\big]\,\dd s\Big\Vert\notag\\
			&\leq\;\int_0^t\norm{\msf{T}_n(s)}_\text{OP}\cdot \norm{\msf{L}_n\Phi_ng_{s,t}-\Phi_n\msf{L}g_{s,t}}\,\dd s\notag \\
			&\overset{\eqref{eq3.1}}{\leq}\;\int_0^t\Big[s_1(n)\norm{g_{s,t}}+2\max\{s_2(n),\,r_1(n)\}\norm{\msf{L}g_{s,t}}\Big]\,\dd s\notag \\ &\qquad+\int_0^t2\max\{s_3(n),\,r_2(n)\}\norm{\msf{L}^2g_{s,t}}\,\dd s
			\notag\\
			&=\!\int_0^t\Big[s_1(n)\norm{\msf{T}(t-s)f}+2\max\{s_2(n),r_1(n)\}\norm{\msf{L}\msf{T}(t - s)f}\Big]\dd s\notag \\ &\qquad+2\int_0^t\max\{s_3(n),r_2(n)\}\norm{\msf{L}^2\msf{T}(t - s)f}\,\dd s
			\notag\\
			&\leq\;2t\Big(s_1(n)\norm{f}+\max\{s_2(n),\,r_1(n)\}\norm{\msf{L}f}\Big)\notag \\
			&\qquad\,+\,2t\max\{s_3(n),\,r_2(n)\}\norm{\msf{L}^2f}\,,\label{last}
		\end{align}
		for every $f \in \mf{D}(\msf{L}^2)$. Note that  hypothesis \ref{(H1)} assures that $\msf{L}_n\Phi_ng_{s,t}$ above is well-defined. Recall that $\Phi_n = \pi_n + \Xi_n$. By the triangle inequality,
		\begin{align*}
			\norm{\msf{T}_n(t)\pi_nf-\pi_n\msf{T}(t)f}\,&\leq\,\norm{\msf{T}_n(t)\Phi_nf-\Phi_n\msf{T}(t)f}\,+\,\norm{\msf{T}_n(t)\Xi_nf}\,+\,\norm{\Xi_n\msf{T}(t)f}.
		\end{align*}
		Invoking hypothesis \ref{(H3)} and applying inequality \eqref{last} we  conclude the proof.
	\end{proof}

\end{document}
